\section{Pointers}

\subsection{Memory Layout and Virtual Address Space}

\begin{definition}
    The operating system provides each process with its own private virtual address space. Memory is organized into specific segments:
    \begin{itemize}
        \item \textbf{Kernel Virtual Memory:} Memory at the top of the address space that is inaccessible to user code.
        \item \textbf{User Stack:} Dynamically created region at runtime for managing function calls, growing downward.
        \item \textbf{Shared Libraries \& Heap:} Memory-mapped region for shared libraries and dynamically allocated memory created via \texttt{malloc}.
        \item \textbf{Read/Write Segment:} Stores initialized and uninitialized global and static variables (\texttt{.data}, \texttt{.bss}) loaded from the executable.
        \item \textbf{Read-Only Segment:} Contains executable code and read-only constants (\texttt{.init}, \texttt{.text}, \texttt{.rodata}) loaded from the executable.
    \end{itemize}
\end{definition}

\begin{remark}
    Set-up code allocates a stack frame upon entering a procedure to store local variables, return information, and temporary space. The space is deallocated by finish code upon return.
\end{remark}

\subsection{Pointer Fundamentals}

\begin{definition}
    A pointer is a variable whose value is the virtual memory address of another variable.
    \begin{itemize}
        \item \textbf{Address-of Operator (\texttt{\&}):} Returns the memory address of a variable.
        \item \textbf{Dereference Operator (\texttt{*}):} Accesses or modifies the memory location referred to by a pointer.
        \item \textbf{Double Pointer:} A variable holding the memory address of another pointer (e.g., \texttt{int **dp}).
        \item \textbf{NULL Pointer:} A pointer assigned to address \texttt{0} (or \texttt{0x00000000}), representing a guaranteed invalid memory location.
    \end{itemize}
\end{definition}

\begin{important}
    Any attempt to dereference a \texttt{NULL} pointer results in a segmentation fault.
\end{important}

\begin{codeexample}
int x = 42;
int *p = &x;   // p holds the address of x
int **dp = &p; // dp holds the address of p
*p = 99;       // dereferences p to assign 99 to x
\end{codeexample}

\subsection{Pointer Operations}

\begin{definition}
    \textbf{Pointer Arithmetic:} Adding or subtracting an integer $n$ to a pointer shifts the target memory address by $n \times \text{sizeof}(\text{type})$ bytes.
\end{definition}

\begin{definition}
    \textbf{Casting:} Type casting a pointer reinterprets the byte-level data at that memory address according to the new type without changing the address itself.
\end{definition}

\begin{codeexample}
int arr[3] = {1, 2, 3};
int *int_ptr = &arr[0];
int_ptr += 1; // moves forward by sizeof(int) bytes to &arr[1]

char *char_ptr = (char *) int_ptr; // casts int* to char*
char_ptr += 1; // moves forward by sizeof(char) byte
\end{codeexample}

\subsection{Function Pointers}

\begin{definition}
    A function pointer stores the executable code address of a function, allowing functions to be passed as arguments or stored inside structures to implement polymorphism and object-oriented patterns.
\end{definition}

\begin{codeexample}
// Pointer to a function taking (int*, char) and returning int
int (*func)(int *, char);

// Example struct containing function pointers (vtable style)
struct file_operations {
ssize_t (*read) (struct file *, char *, size_t, loff_t *);
ssize_t (*write) (struct file *, const char *, size_t, loff_t *);
};
\end{codeexample}
